\chapter{Uvod}

Transport robe teškim teretnim vozilima podleže ograničenjima koja ne postoje kod putničkih automobila: visina vozila mora biti manja od visine svakog nadvožnjaka i tunela na ruti, ukupna masa i osovinsko opterećenje moraju biti u skladu sa nosivošću mostova, širina i dužina vozila ograničavaju kretanje kroz uža naselja i oštre krivine, a vozila koja prevoze opasan teret moraju zaobilaziti određene deonice (tunele, naseljena mesta) po propisu. Standardne aplikacije za navigaciju namenjene putničkim vozilima (Google Maps, Waze i slične) te podatke ne uzimaju u obzir, pa ruta koju predlažu vozaču teretnog vozila može biti fizički neprohodna ili čak zakonski nedozvoljena za dato vozilo. Sa druge strane, postojeća komercijalna rešenja specijalizovana za teretni saobraćaj (npr. PTV navigator, Google Maps za dostavu, TomTom Truck) su zatvorena, plaćena i ne omogućavaju uvid u to \emph{zašto} je određena ruta izabrana, što otežava vozaču da proceni da li je predlog razuman.

Ovaj rad opisuje razvoj sistema za rutiranje teretnih vozila koji rešava taj problem korišćenjem otvorenih geografskih podataka projekta OpenStreetMap (OSM) i otvorenog routing engine-a Valhalla, koji već ima ugrađen model cene rute specifičan za teretna vozila (\emph{truck costing}), sposoban da isključi deonice koje ne zadovoljavaju ograničenja visine, mase, širine ili prevoza opasnog tereta datog vozila. Na tu osnovu je dodat sopstveni sloj koji:

\begin{enumerate}
\item rangira alternativne rute koje Valhalla vrati prema riziku i preferencama vozača (stil vožnje, osetljivost tereta, potrošnja goriva, favorizovane lokacije na mapi), pošto sam Valhalla to ne radi;
\item vozaču objašnjava \textbf{zašto} je predložena ruta odstupila od geometrijski najkraćeg puta, umesto da mu samo prikaže krajnji rezultat;
\item omogućava, kao samostalnu algoritamsku celinu odvojenu od produkcionog puta, sopstvenu implementaciju Dijkstra i A* pretrage najkraćeg puta nad ograničenim podgrafom putne mreže, radi demonstracije i evaluacije principa na kojima Valhalla interno počiva.
\end{enumerate}

Sistem je zamišljen ne samo kao alat za planiranje jedne rute, već kao kompletna, iako po obimu ograničena, distribuirana aplikacija za rad voznog parka: pored planiranja i praćenja pojedinačne ture, uveden je i model dispečera koji upravlja više vozača i vozila, mehanizam ponude i prihvatanja ture između dispečera i vozača, chat u realnom vremenu, praćenje pozicije vozila uživo preko WebSocket veze, dnevnik događaja tokom vožnje i jednostavan mehanizam predloga pauze vozača zasnovan na proteklom vremenu vožnje. Sistem se sastoji od Go backend servisa, PostgreSQL/PostGIS baze podataka, RabbitMQ posrednika poruka, Valhalla routing engine-a i Flutter mobilne aplikacije, a svi servisi su kontejnerizovani i pokreću se preko \texttt{docker-compose}.

\section{Motivacija}

Motivacija za rad proizlazi iz konkretnog, ponovljivog problema: standardna navigacija ne razlikuje teretno vozilo od putničkog automobila, dok specijalizovana komercijalna rešenja skrivaju logiku odabira rute od korisnika i po pravilu zahtevaju plaćenu licencu i vlasnički skup podataka o putnoj mreži. Cilj ovog rada je da pokaže da se prihvatljivo tačan i \textbf{objašnjiv} sistem za rutiranje teretnih vozila može izgraditi nad potpuno otvorenim podacima (OSM) i otvorenim routing engine-om (Valhalla), uz sopstveni algoritamski doprinos koji tim alatima nedostaje: procenu rizika rute specifičnu za teretni saobraćaj i automatsko objašnjenje odstupanja rute od očekivanog, najkraćeg puta.

\section{Ciljevi i obim rada}

Ciljevi rada su:

\begin{itemize}
\item proučiti postojeća rešenja za rutiranje teretnih vozila i teorijske osnove distribuiranih sistema i algoritama pretrage najkraćeg puta u grafu;
\item pripremiti geografske podatke za teritoriju Republike Srbije iz OSM ekstrakta, sa oznakama relevantnim za teretni saobraćaj;
\item implementirati backend servis koji generiše, rangira i objašnjava rutu za dato vozilo, koristeći Valhallu kao routing engine;
\item implementirati, kao odvojenu algoritamsku celinu, sopstvenu Dijkstra/A* pretragu nad ograničenim podgrafom, radi evaluacije principa rutiranja bez zavisnosti od Valhalla-e;
\item implementirati distribuiranu komunikaciju sistema (REST, asinhrona obrada, WebSocket) i mobilnu aplikaciju za vozače i dispečere;
\item testirati implementirano rešenje i evaluirati doprinos sopstvene funkcije rizika u odnosu na "sirov" izbor rute.
\end{itemize}

Rad se svesno ne bavi punom implementacijom evropske AETR regulative o radnom vremenu i pauzama vozača, elevacionim (nagibnim) rizikom rute, offline režimom rada mobilne aplikacije, podrškom za više država niti bezbednosnim hardeningom na produkcionom nivou — ta ograničenja su namerna i obrazložena u poglavlju 10.

\section{Struktura rada}

Poglavlje 2 daje pregled postojećih komercijalnih i otvorenih rešenja za rutiranje teretnih vozila i teorijske osnove (distribuirani sistemi, algoritmi pretrage najkraćeg puta, OSM podatkovni model).\\Poglavlje 3 opisuje korišćene tehnologije.\\Poglavlje 4 opisuje arhitekturu sistema. \\Poglavlje 5 opisuje pripremu geografskih podataka.\\Poglavlje 6 detaljno opisuje algoritamski deo rada — Valhalla truck costing, sopstvenu funkciju rizika, mehanizam objašnjenja rute i sopstvenu Dijkstra/A* implementaciju. \\Poglavlje 7 opisuje backend servis i model podataka. \\Poglavlje 8 opisuje mobilnu aplikaciju. \\Poglavlje 9 opisuje testiranje i evaluaciju. \\Poglavlje 10 sadrži zaključak i pravce budućeg rada.

\clearpage
